\section{Introduction}
\label{sec:introduccion-ch}

The Continuum Hypothesis asks whether there exists a set of real numbers
whose cardinality is strictly greater than that of the natural numbers and
smaller than that of the complete real line. This work gives a negative answer
to that existence question, and hence an affirmative answer to the hypothesis,
in a genealogical closure defined through the operations of Triadic Modular
Holography (HMT). Its contribution links the construction of the real line by
compatible prolongations to the cardinal comparison of all its internal
subsets.

The HMT formal core comprises three objects developed from their definitions.
All Possible Paths (APP) equips the support \((\mathbb Z/9\mathbb Z)^2\)
with sum and product evaluations, retaining their residues and quotients.
Translations by \((\pm1,0)\) and \((0,\pm1)\) determine its Cayley graph;
square cells and the identification of opposite boundaries give it a
two-dimensional toroidal realization. TRIT denotes the oriented ternary
signature \(\{-1,0,+1\}\), which distinguishes transition regimes. The
Topological Phase Kernel (TPK) composes selection, transport, and memory to
prolong states. These distinct functions act jointly: the support locates
paths, the evaluations determine their arithmetic content, and the dynamics
preserves the history of each transition.

The argument follows this composition from finite operations to the domain in
which the cardinal comparison is performed:
\[
 \mathrm{APP}\longrightarrow\mathrm{TRIT}\longrightarrow\mathrm{TPK}
 \longrightarrow X_\infty^{\mathrm{enr}}
 \longrightarrow\mathcal C_{\mathrm{cont}}^{\mathrm{disc}}
 \longrightarrow m_\infty
 \longrightarrow\mathfrak G_{\mathrm{HMT}}(h,m_\infty).
\]
The composition produces enriched states, compatible prolongations, complete
histories, and only thereafter an Archimedean realization. The real line is an
output of the construction, not its initial universe; the cardinal question
is decided over the entire domain produced by that genealogy.

The central question is whether the line obtained by Archimedean projection admits,
within its complete genealogical closure, a cardinal strictly intermediate
between \(\aleph_0\) and the cardinal of the line itself. The theorem is
formulated in \(\mathfrak G_{\mathrm{HMT}}(h,m_\infty)\), the closure generated
by the complete system. In that domain it quantifies over
\(\mathcal P^{\mathfrak G}(\mathbb R_{\mathrm{HMT}}^{\mathfrak G})\) and
establishes the complete cardinal dichotomy.

Theorem~\ref{thm:decision-continuo-hmt}, proved in
Section~\ref{sec:cierre-cardinal-nativo}, establishes the main result. In the
cardinal comparisons, the real line and the power sets are interpreted in the
genealogical closure
\(\mathfrak G=\mathfrak G_{\mathrm{HMT}}(h,m_\infty)\). This domain and its
quantifiers are constructed explicitly before the comparison is performed.
The result is as follows. For every realized history \(m_\infty\) produced by
the complete system associated with
\[
 \mathrm{APP}\longrightarrow\mathrm{TRIT}\longrightarrow\mathrm{TPK}
 \longrightarrow X_\infty^{\mathrm{enr}}
 \longrightarrow\mathcal C_{\mathrm{cont}}^{\mathrm{disc}},
\]
whose restrictions are compatible and whose prolongations are nonempty,
exhaustive, and surviving at every horizon, the genealogical extension
\(\mathfrak G_{\mathrm{HMT}}(h,m_\infty)\) generates the integral real line,
quantifies over its entire power set, and satisfies
\[
 \mathfrak G_{\mathrm{HMT}}(h,m_\infty)
 \models 2^{\aleph_0}=\aleph_1.
\]
In particular, for every
\(A\in\mathcal P^{\mathfrak G}(\mathbb R_{\mathrm{HMT}}^{\mathfrak G})\),
\[
 |A|\leq\aleph_0
 \qquad\text{or}\qquad
 |A|=|\mathbb R_{\mathrm{HMT}}^{\mathfrak G}|.
\]
Consequently, HMT proves the Continuum Hypothesis affirmatively for the
integral continuum that it defines and generates. The formal statement and
its proof are given once in the main theorem.

The essential difference from a purely extensional presentation lies in the
identity of the mathematical object. An isolated real coordinate preserves a
value; an HMT state also preserves the operations that produced it, the
arithmetic sheet, the tritic orientation, the residues and quotients traversed,
the carry, the boundary, the route, and the accumulated memory. The
Archimedean projection neither defines nor exhausts that structure: it
determines one of its coordinates. The second projection preserves the boundary
incidence that, in downstream developments, reaches
Paley--Witt--Golay--Leech, \(V^\natural\), and the Monster. Both realizations
arise from the same state, and neither is introduced to select the other
retrospectively.

\subsection{Construction of the continuum through compatible prolongations}

APP arranges the sum and product sheets over the same support and preserves
the difference between their additive and multiplicative evaluations. TRIT incorporates orientation and
regime. TPK selects, transports, memorizes, and prolongs the admissible
transitions. The composition does not terminate in a finite state: it defines
a system of restrictions
\[
 \rho_{N+1,N}:X_{N+1}\longrightarrow X_N
\]
and nonempty extensions
\[
 \operatorname{Ext}_N(x)\subseteq\rho_{N+1,N}^{-1}(x).
\]
Compatibility at every depth produces the inverse limit
\[
 X_\infty^{\mathrm{enr}}
 =\varprojlim_N(X_N,\rho_{N+1,N}).
\]
The quantifier is \(\forall N\); no execution to one thousand, ten thousand,
or one million digits constitutes the bound of the theorem.

Each successor block subdivides a cylinder into \(729=3^6\) disjoint,
exhaustive refinement subcylinders of decreasing diameter. This number
describes the local branching of the refinement; it is not identified with
the cardinal of the continuum. The complete block tape determines a history,
and its arithmetic projection determines a unique real point. Conversely, the
canonical expansion of each real reconstructs a tape, and the coverage
property lifts it to a compatible HMT history. Thus the real line is a
surjective image of the constructed tree, never an external set used to
choose its branches.

Measurement enters as a coupling. On the domain under consideration, the
fiber-separation condition is imposed first: the enriched reader must
distinguish two different states that yield the same value under Archimedean evaluation. Under that
condition, formalized in Proposition~\ref{prop:ch-recuperabilidad-fibra}, the
instrument distinguishes a compatible continuation, adds the result and the
back-action to the genealogical register, and restricts future histories to
those that prolong the new state. The realized history \(m_\infty\) is formed
through these updates. The inverse transformation then reconstructs the
interior from the memory transported to the boundary, so that the update does
not amount to a loss of information. Proposition~\ref{prop:ch-naturalidad-medida}
shows when changes of representation also preserve the measure of histories.

\subsection{Genealogical semantic closure}

Let \(h\) be the effective code of the graphs of APP, TRIT, and TPK, the
prolongations, the restrictions, and the two projections. The pair
\((h,m_\infty)\) is primitively coded by a single real
\(u_{h,m}=h\oplus m_\infty\). Define
\[
 \begin{aligned}
 \mathfrak G_0(h,m_\infty)
   &:=\operatorname{TC}\{\omega,u_{h,m},h,m_\infty\},\\
 \mathfrak G_{\beta+1}(h,m_\infty)
   &:=\operatorname{Def}_{\Sigma_{\mathrm{HMT}}}
        \bigl(\mathfrak G_\beta(h,m_\infty)\bigr),\\
 \mathfrak G_\lambda(h,m_\infty)
   &:=\bigcup_{\beta<\lambda}\mathfrak G_\beta(h,m_\infty),\\
 \mathfrak G_{\mathrm{HMT}}(h,m_\infty)
   &:=\bigcup_{\beta\in\operatorname{Ord}}
        \mathfrak G_\beta(h,m_\infty).
 \end{aligned}
\]
This closure is formulated through relative definability and condensation in
the precise sense of constructible hierarchies~\cite{Jech2003}. The Continuum
Hypothesis is not introduced among the conditions defining the closure; it is
obtained by the proof developed in this article. Neither an enumeration of the
real line nor a target bijection is introduced as data. Its successor operation
forms the subsets definable by finite formulas of the HMT language and
parameters already produced. Equivalently, every admitted object has a
derivation tree, and every admissible tree determines an object.

The resulting hierarchy is transitive, contains its ordinals, satisfies the
axioms employed in the cardinal comparison, and has a global well-order
determined by the first level of birth and the least derivation name. The
notation
\[
 \mathfrak G_{\mathrm{HMT}}(h,m_\infty)
 =L[u_{h,m}]=L[h,m_\infty]
\]
is obtained later as a representation theorem. It is not the engine of APP,
TRIT, or TPK, it does not select \(m_\infty\), and it does not replace the
material construction of the continuum.

\subsection{Cardinal comparison between the HMT real line and the first uncountable ordinal of the closure}

The cardinal proof is divided into two directions independent of the equality
sought. First, genealogical condensation proves that every generated real
appears before \(\omega_1^{\mathfrak G}\). To retain the convention used in
the proof, let \(\iota(r)\) be the code of the rational cut of \(r\), let
\(\beta(r)\) be the least ordinal such that
\(\iota(r)\in\mathfrak G_{\beta(r)+1}\), and let \(n(r)\) be the least index
of this cut in the enumeration \(e_{\beta(r)+1}\). Define
\[
 j_{\mathrm{HMT}}(r)=\omega\beta(r)+n(r),
 \qquad
 j_{\mathrm{HMT}}:
 \mathbb R_{\mathrm{HMT}}^{\mathfrak G}
 \hookrightarrow\omega_1^{\mathfrak G}.
\]

In the opposite direction, every ordinal
\(\beta<\omega_1^{\mathfrak G}\) has a real code \(w_\beta\) for a countable
well-order of type \(\beta\). The codes are chosen by means of the
genealogical well-order already constructed. The ternary embedding
\[
 E(w)=\sum_{n\geq0}\frac{2\,\mathbf 1_w(n)}{3^{2n+2}}
\]
is injective, and the uniform cylinder lift proves that \(E(w_\beta)\) belongs
to the real line obtained by the HMT Archimedean projection. Therefore,
\[
 k_{\mathrm{HMT}}:
 \omega_1^{\mathfrak G}
 \hookrightarrow\mathbb R_{\mathrm{HMT}}^{\mathfrak G}.
\]
Cantor--Bernstein closes the cardinal equality. The proved identification
\[
 \mathbb R_{\mathrm{HMT}}^{\mathfrak G}
 \approx\mathcal P^{\mathfrak G}(\omega)
\]
turns it exactly into \(2^{\aleph_0}=\aleph_1\). Here \(A\approx B\) denotes
equinumerosity; it does not assert a field isomorphism between a real line and
a power set.

\clearpage
\subsection{Logical structure of the proof}

The proof is organized into four logical steps. First, the domain of histories,
its inverse limit, and the exhaustive coverage of cylinders are constructed.
Second, the semantic closure is formed and condensation is applied to its
reals. Third, without using the equality sought, the two injections between
the HMT real line and \(\omega_1^{\mathfrak G}\) are defined. Fourth,
Cantor--Bernstein yields the equality, and the cofinality argument establishes
the dichotomy for every internal subset. In this sequence, \(729\) describes
local branching; coverage relates histories and values; condensation controls
the level at which internal reals appear. The constants are correlated
outputs of the same state: their conventional digits do not select the
cardinal comparison, while their regions, routes, and memories remain in the
enriched object on which the real line is constructed. The subsequent
identification with \(L[h,m_\infty]\) represents the closure already
constructed.\looseness=-1

The order of exposition first reproduces the formal core shared by the series,
the prolongation, the joint construction of the continuum, and the coinductive
generation; it then presents the integral semantic closure and the cardinal
proof. Only thereafter does it develop the reduced phase-and-memory core, the
integral dodecaphase seal \(K\), the complete dodecaphase route to
\(\alpha_{\mathrm{HMT}}\), the finite analytic witness of order nine, the
correlated analytic representation with recurrence, convergence, and common
cylinders, and the exceptional corridor. These coordinates share an originating
state but have different codomains and functions; jointly, they exhibit the
complete dual realization and preserve the evaluation and projection routes, incidence,
boundary, and memory. No target value selects the cardinal injections. The
comparison with Cantor, Gödel, and Cohen is introduced at the end of the
construction~\cite{Cantor1878,Hilbert1902,Godel1938,Cohen1963,Cohen1964} and
does not govern the definition of the HMT object.\looseness=-1

\enlargethispage{2\baselineskip}
Every arrow used by the cardinal closure is defined in the body of the article
by its domain, codomain, action, and compatibility with refinement: the three
oriented loops of the nonadic return, their grouping into six-symbol blocks,
the \(54\)-step block extension, the construction of cylinders, and the uniform
lift of their branches. The nine transitions are constructed explicitly and
provide the uniform lift used in the cardinal comparison. The reproducible
calculations accompany, but do not replace, this proof.\looseness=-1

% BEGIN R27_FRAMING_INTRO_EN
The transition from seeds to compatible histories contains finite
readings distinguished by the information they retain. The additive and
multiplicative signatures determine the \(42\) possible blocks of one
decimal window, with explicit seeds proving reachability. Six-window
emissions establish a junction cycle between regional halves. Temporal
transport retains phase, orientation, and memory in addition to those
incidences. A pair of dodecaphasic records with the same histogram and
different correlation proves why the archive requires sample order.
These proofs specify the finite antecedent of prolongation; coverage,
semantic closure, and cardinal injections are developed in their own
domains.
% END R27_FRAMING_INTRO_EN

% BEGIN R28_FRAMING_INTRO_EN
The finite foundational part includes the emission multigraph, its
phase quotient, and five minimal eight-edge closures. Each closure
jointly covers APP, propagation, autoscale, and one closure region.
The oriented witnesses, finite search, and parity of memory readers
specify the data retained before the prolongations and cardinal
readings developed in this work.
% END R28_FRAMING_INTRO_EN
